% ============================================================
%  陳述書テンプレート  (kian-v08 LuaLaTeX 版)
%  コンパイル: latexmk -r latexmkrc-lua template-chinjutsu-lua.tex
%              または VS Code（LaTeX Workshop）でビルド
%  プレースホルダー: 【...】を実際の内容に置き換えてください
%  異体字: \ivs{文字} で IPAmj 明朝字形に切り替え（例: \ivs{角󠄁}）
%  ※ \SecA{} 等の {} 内に % を書くとコメントになりエラーが出ます
%  見出し番号: \SecZ=番号なし \SecA=第１ \SecB=１ \SecC=⑴ \SecD=ア \SecE=(ア)
%  ※ 見出し（\SecA{} 等）は本文なしで連続させると番号表示が崩れるため、
%    各見出しの直後に【ここに記載】のプレースホルダー本文を入れています
% ============================================================

\documentclass{ltjsarticle}
\usepackage{kian-lua}
% \usepackage{kian-ivs} %IVS 異体字（\ivs{文字}）
% \usepackage{kian-otf} %囲み文字拡張
% \usepackage{luatexja-ruby} %ルビ機能

\Idcs{
  [o]{R８(　)　　:【事件名】\\原告:【氏/名】:\\被告:【氏/名】}
}

%  陳述書の表題: 陳述者の氏名を括弧内に記す
%  氏名に異体字がある場合: 陳述書（\ivs{角󠄁}田　一郎）
\Tit{陳述書（【陳述者氏名】）}
\Dat{令和８年　　月　　日}
\Recs{
  [n]{:【裁判所名】【支部名】支部:御中}
}

% ============================================================

\begin{document}

\maketitle

% ---------------------------------------------------------- %
%  自己紹介
%  ※ 陳述書は一人称（「私は」）で記載するのが原則
% ---------------------------------------------------------- %

\SecA{自己紹介}



% ---------------------------------------------------------- %
%  本件との関係
% ---------------------------------------------------------- %

\SecA{本件との関係}



% ---------------------------------------------------------- %
%  本件の経緯
% ---------------------------------------------------------- %

\SecA{本件の経緯}

\SecB{}

【ここに記載】

\SecB{}

【ここに記載】

\SecB{}

【ここに記載】

% ---------------------------------------------------------- %
%  結語
% ---------------------------------------------------------- %

\SecA{結語}


\vspace{3\zw}

% ---------------------------------------------------------- %
%  陳述者情報（押印欄）
%  氏名に異体字がある場合: \ivs{角󠄁}田　一郎
% ---------------------------------------------------------- %

\noindent\Jus[4]{住　所}　【陳述者の住所】

\noindent\Jus[4]{氏　名}　【陳述者の氏名】　　　　印

\Ijo[2\zw]

\end{document}
